(1) 因 \(C=A+a\mathbf e_1\)，且 \(\dot{\mathbf e}_1=\omega\mathbf e_2\)，
\[
\mathbf v_C=u\mathbf e_1+a\omega\mathbf e_2,\qquad
\mathbf a_C=(\dot u-a\omega^2)\mathbf e_1+(u\omega+a\dot\omega)\mathbf e_2.
\]
于是 \(T=\tfrac12m u^2+\tfrac12(I+ma^2)\omega^2\)。设反力 \(\mathbf F=\Lambda\mathbf e_2\)，纵向牛顿方程给 \(\dot u=a\omega^2\)；横向给 \(\Lambda=m(u\omega+a\dot\omega)\)。反力对质心的力矩是 \(-a\Lambda\)，故 \(I\dot\omega=-a\Lambda\)。联立可得
\[
\boxed{\dot u=a\omega^2,\qquad J\dot\omega=-ma\,u\omega,\qquad
\Lambda=\frac{mI}{J}u\omega},\quad J=I+ma^2.
\]
该条件约束瞬时速度方向，不能通过对一个固定位置方程求导来代替。

(2) 能量导数
\[
\dot E=mu\dot u+J\omega\dot\omega
=ma\,u\omega^2-ma\,u\omega^2=0.
\]
定义 \(v=\sqrt{2E/m}\)、\(\Omega=v\sqrt{m/J}\)、\(\kappa=mav/J\)。能量椭圆为 \(u^2/v^2+\omega^2/\Omega^2=1\)。沿非平凡轨道 \(\omega\) 符号 \(\sigma=\pm1\) 不变；分离变量积分 \(\dot u=(ma/J)(v^2-u^2)\) 给
\[
\boxed{u=v\tanh[\kappa(t-t_0)],\quad
\omega=\sigma\Omega\,\operatorname{sech}[\kappa(t-t_0)]}.
\]
直接求导：\(\dot u=v\kappa\operatorname{sech}^2=a\omega^2\)，\(\dot\omega=-\kappa\tanh[\kappa(t-t_0)]\omega=-(ma/J)u\omega\)。若 \(\omega(0)=0\)，则 \(\omega\equiv0,\ u\equiv\pm v\) 是特殊直线解，不能用有限 \(t_0\) 表示。

(3) 因 \(\int\operatorname{sech}y\,dy=2\arctan(e^y)\)，积分得到
\[
\boxed{\theta(t)=\theta_-+2\sigma\frac{\sqrt{J/m}}a
\arctan(e^{\kappa(t-t_0)})},\qquad
\boxed{\theta_+-\theta_-=\sigma\pi\frac{\sqrt{J/m}}a}.
\]
在 \(u=v\) 附近，\(\delta\dot\omega=-\kappa\delta\omega\)，故角扰动衰减；在 \(u=-v\) 附近，\(\delta\dot\omega=+\kappa\delta\omega\)，故不稳定。非直线轨道由过去的 \(-v\) 连接到未来的 \(+v\)。当 \(a\to0\)，\(\kappa\to0\) 而总转角 \(\propto1/a\)，两个极限不交换；恰好 \(a=0\) 时 \(u,\omega\) 均为常数。

(4) 设 \(s=t-t_0\)、\(\omega_0=\sigma\Omega\)、\(\mathbf e_{j0}=\mathbf e_j(t_0)\)。由方程有 \(u(t_0)=0\)、\(\dot u(t_0)=a\omega_0^2\)、\(\ddot u(t_0)=0\)、\(\dot\omega(t_0)=0\)。展开 \(\dot{\mathbf r}_A=u\mathbf e_1(\theta)\) 并积分：
\[
\boxed{\mathbf r_A(t_0+s)-\mathbf r_A(t_0)
=\frac{a\omega_0^2}{2}s^2\mathbf e_{10}
+\frac{a\omega_0^3}{3}s^3\mathbf e_{20}+O(s^4)}.
\]
故局部横坐标 \(\propto s^2\)、纵坐标 \(\propto s^3\)，形成半三次尖点；此时刀刃瞬时静止但质心仍可旋转。反力 \(\Lambda=(mI/J)u\omega\) 在尖点为零，且理想约束功率 \(\mathbf F\cdot\mathbf v_A=\Lambda\mathbf e_2\cdot u\mathbf e_1=0\) 恒成立。检查 \([\kappa]=T^{-1}\)、\([\Lambda]=MLT^{-2}\) 和无量纲转角均正确。
